% =========================================================
% SECCIÓN 2.4 - PROTOCOLO DE MEDICIÓN Y CONVERSIÓN A FLUJO MÁSICO
% =========================================================

La cuantificación de flujos gaseosos en sistemas de bioconversión con sustrato sólido exige un compromiso entre trazabilidad de masa y mínima perturbación del sistema. El protocolo de ventanas cerradas (\textit{closed static chambers}) se adopta como solución de campo ante la imposibilidad de usar cámaras de flujo continuo sin alterar la oxigenación larval \parencite{rossiEstimatingDynamicsGreenhouse2024a, parodiBioconversionEfficienciesGreenhouse2020a}. En esta sección se describe el ciclo de muestreo y se deriva formalmente la conversión de concentración (\si{ppm}) a flujo másico instantáneo (\si{g/min}), puente metrológico entre el sensor NDIR y el balance de carbono del modelo DEB.

% ---------------------------------------------------------
\subsection{Diseño del ciclo de muestreo}
\label{subsec:cap2_ciclo_muestreo}

El protocolo se estructuró en eventos discretos espaciados cada 2 días. En cada evento se cerraba herméticamente el reactor (tapa de plástico con sellado por presión) y se iniciaba el registro continuo de acumulación gaseosa hasta la saturación del instrumento (\num{5000}~ppm para \ce{CO2} o \SI{10}{\percent} para \ce{CH4}). Las lecturas se tomaron cada \SI{10}{\second} en un microcontrolador ESP32, descartando los primeros \SI{2}{min} para homogeneización del \emph{headspace}. La duración efectiva de cada ventana varió biológicamente: estadios tempranos (L3--L4) requirieron tiempos prolongados, mientras que L5 acortó el intervalo por respiración elevada. Esta variabilidad captura la dinámica real sin truncar artificialmente picos de emisión.

El reactor consistió en un recipiente rígido de \SI{0.8}{L} con \num{700} larvas aproximadamente y \SI{250}{g} de sustrato húmedo, dejando un volumen de cabeza libre $V_{\text{head}} = \SI{0.35}{L}$ ($\pm$\,\SI{0.05}{L}). El número de eventos válidos por réplica osciló entre 6 y 7 a lo largo de un ciclo de 12--\SI{15}{\day}.

% ---------------------------------------------------------
\subsection{Fundamento teórico de la conversión ppm--flujo másico}
\label{subsec:cap2_fundamento_teorico_conversion}

La señal cruda $C(t)$ (\si{ppm}) se convierte en flujo másico mediante tres principios físico-químicos acoplados: ecuación de estado del gas ideal, ley de Dalton de presiones parciales y balance de masa transitorio en un volumen de control de paredes rígidas.

% - - - - - - - - - - - - - - - - - - - - - - - - - - - - -
\subsubsection{Hipótesis de gas ideal y densidad molar}

En las condiciones del ensayo ($P_{\text{atm}} \approx \SI{101.3}{kPa}$, $T_{\text{air}} \approx \SI{303}{K}$) la mezcla gaseosa (aire + trazas de \ce{CO2}) se comporta como gas ideal con error despreciable ($<<$\SI{0.3}{\percent}) \parencite{atkinsPhysicalChemistry2014}. La densidad molar del \ce{CO2} es:

\begin{equation}
\rho_{\ce{CO2}} = \frac{P_{\text{atm}} \cdot M_{\ce{CO2}}}{R \cdot T_{\text{air}}}
\label{eq:densidad_co2}
\end{equation}

donde $M_{\ce{CO2}} = \SI{44.01}{g/mol}$ y $R = \SI{8.314}{kPa\cdot L/(mol\cdot K)}$. En Palmira resulta $\rho_{\ce{CO2}} \approx \SI{1.80}{g/L}$.

% - - - - - - - - - - - - - - - - - - - - - - - - - - - - -
\subsubsection{Fracción volumétrica y masa acumulada}

El sensor NDIR reporta concentración volumétrica $C_{\ce{CO2}}$ en ppm, definida por la IPCC como la razón molar respecto al total de la mezcla multiplicada por $10^6$ \parencite{ipccGuidelinesNationalGreenhouse2006}:

\begin{equation}
C_{\ce{CO2}}(t) = \frac{n_{\ce{CO2}}(t)}{n_{\text{total}}} \times 10^{6}
\quad\Longrightarrow\quad
y_{\ce{CO2}}(t) = C_{\ce{CO2}}(t) \times 10^{-6}
\label{eq:definicion_ppm}
\end{equation}

donde $y_{\ce{CO2}}$ es la fracción molar, igual a la fracción volumétrica para gases ideales. La masa acumulada en el \emph{headspace} es:

\begin{equation}
m_{\ce{CO2}}(t) = y_{\ce{CO2}}(t) \cdot \rho_{\ce{CO2}} \cdot V_{\text{head}}
= \frac{M_{\ce{CO2}} \cdot P_{\text{atm}} \cdot V_{\text{head}}}{R \cdot T_{\text{air}}} \cdot y_{\ce{CO2}}(t)
\label{eq:masa_teorica}
\end{equation}

% - - - - - - - - - - - - - - - - - - - - - - - - - - - - -
\subsubsection{Balance de masa transitorio y flujo másico}

El reactor batch cerrado es un volumen de control sin entrada de masa gaseosa durante la medición. La única fuente de \ce{CO2} es la respiración biológica del sistema larva-sustrato. Derivando la \cref{eq:masa_teorica} respecto al tiempo (con $P_{\text{atm}}$, $V_{\text{head}}$ y $T_{\text{air}}$ constantes en la ventana corta de medición):

\begin{equation}
\dot{m}_{\ce{CO2}}(t) 
= \frac{\mathrm{d} m_{\ce{CO2}}}{\mathrm{d} t}
= \underbrace{\frac{M_{\ce{CO2}} \cdot P_{\text{atm}} \cdot V_{\text{head}}}{R \cdot T_{\text{air}}} \times 10^{-6}}_{\displaystyle\kappa}
\cdot \frac{\mathrm{d} C_{\ce{CO2}}}{\mathrm{d} t}
\label{eq:flujo_masico_derivado}
\end{equation}

El factor $\kappa$ (\si{g\cdot min/ppm}) actúa como coeficiente de proporcionalidad entre la pendiente de concentración y el flujo másico instantáneo. En la práctica, la derivada se aproxima por regresión lineal sobre la serie $C(t)$ durante la ventana de acumulación, descartando los puntos iniciales de equilibración ($t < \SI{2}{min}$). La pendiente $\Delta C/\Delta t$ (\si{ppm/min}) se sustituye en la \cref{eq:flujo_masico_derivado} para obtener el flujo promedio de la ventana.

La linealización es válida si la tasa de producción biológica varía lentamente respecto al tiempo de integración (hipótesis de \emph{quasi-estacionariedad}). En este reactor, las ventanas de saturación duran entre \SI{10}{min} y \SI{60}{min}, intervalo en que el metabolismo larval no experimenta cambios abruptos de escala fisiológica \parencite{eriksenDynamicModellingFeed2022}. El coeficiente de determinación $R^2 > 0.95$ garantiza que el error de linealización es inferior al \SI{2}{\percent}.

% - - - - - - - - - - - - - - - - - - - - - - - - - - - - -
\subsubsection{Corrección por fugas}

El sellado por presión no es hermético perfecto. La fuga sigue una cinética de primer orden proporcional al gradiente de concentración \parencite{hutchinsonMethodsSoilAnalysis1992}:

\begin{equation}
\frac{\mathrm{d} C_{\text{int}}}{\mathrm{d} t}\bigg|_{\text{fuga}} 
= -\lambda \bigl(C_{\text{int}} - C_{\text{ext}}\bigr)
\label{eq:decaimiento_fuga}
\end{equation}

con $C_{\text{ext}} \approx \SI{400}{ppm}$. La constante $\lambda \approx \SI{0.003}{min^{-1}}$ se determinó experimentalmente mediante ensayos de decaimiento con inyección de \ce{CO2} patrón en reactores sin sustrato. El flujo corregido resulta:

\begin{equation}
\dot{m}_{\text{corr}} = \dot{m}_{\text{obs}} + \lambda \cdot m_{\ce{CO2}}(t)
\label{eq:flujo_corregido}
\end{equation}

Esta corrección fue menor al \SI{5}{\percent} del flujo total en todas las réplicas; se conserva en la formulación para garantizar trazabilidad metrológica.

% - - - - - - - - - - - - - - - - - - - - - - - - - - - - -
\subsubsection{Propagación de incertidumbre}

La incertidumbre del flujo másico se propaga conforme al GUM (JCGM 100:2008) \parencite{jcgmEvaluationMeasurementData2008} a partir de fuentes tipo A (repetibilidad de la pendiente $\Delta C/\Delta t$) y tipo B (resolución del sensor MH-410D: \SI{50}{ppm}; precisión de $V_{\text{head}}$: $\pm$\,\SI{0.05}{L}; variabilidad de $T_{\text{air}}$: $\pm$\,\SI{2}{\celsius}; presión: $\pm$\,\SI{0.5}{kPa}). La incertidumbre expandida $U = k \cdot u_c$ con $k=2$ ($\approx 95\%$ de confianza) resulta del orden de $\pm$\,\SI{12}{\percent} para $\dot{m}_{\ce{CO2}}$, valor aceptable para un sistema de bajo costo en contexto agroindustrial \parencite{biagiDevelopmentMachineLearningbased2024a, parodiBioconversionEfficienciesGreenhouse2020a}.

El mismo procedimiento se aplica al \ce{CH4} con $M_{\ce{CH4}} = \SI{16.04}{g/mol}$. Sin embargo, dado que sus concentraciones fueron frecuentemente inferiores al umbral de detección operativo del MH-441D en condiciones predominantemente aeróbicas, los flujos de \ce{CH4} se reportaron como $<$LOQ salvo en eventos de anaerobiosis detectados (Capítulo 5).